\documentclass{ltjsbook}
\usepackage{docmute} 
\usepackage{../../myStyle} 

\begin{document}
\chapter{テスト理論}
\label{H02_TestTheory}
前回は態度尺度の作り方を見てきました。態度は個人の中にある比較的安定した傾向性であり，刺激ー反応を媒介するものと考えれています。これはいわゆるS-R図式，つまり刺激と反応のセットから，ブラックボックスの性質を表現しようという\keyterm{機能主義}{きのうしゅぎ}{functionalism}の考え方そのものです。\textit{funciton}は関数という意味もありますから，\keytermL{関数主義}{かんすうしゅぎ}と言い換えてもいいかもしれません。

直接は目に見えない，潜在的な状態を測定しようという発想は\textbf{テスト}と同じ発想です。皆さんが幾度となく受けてきたあのテスト，学力テストです。テストは「正答数が多ければ点数が高く，正答数が低ければ点数が低い」ものであり，幾つもの問題を被検者に投げかけて，「正答する」という一貫した傾向が見られれば，その人の中に「知識」「学力」「能力」といった性質があるだろうと考える，というものだからです。

心の状態(態度)も一定の刺激に対する反応パターンから類推するものですので，テスト理論の考え方は態度測定の話と非常に密接な関係があります。そこでここでは，測るということについて理論的な裏付けを積み重ねてきたテスト理論を概観して，心の測定の参考にしてみましょう。

\subsection{古典的テスト理論}
心理学における測定は，物理学などハードサイエンスと違って誤差が大きいことを見越して考えなければなりません。知能検査，性格特性などを測定するために考え出された古典的テスト理論のモデルは，
\[X=t+e\]
で表されます。すなわち，テストの点数$X$は真のスコア$t$と誤差$e$に分割できるというものです。
非常に単純なモデルですが，測定したものには誤差がついているという考え方，言い換えると目に見えるものだけが真実ではないという考え方がしめされています。この考え方は，ソフトサイエンスの領域においては重要なことです。

またこの古典的テスト理論から，いくつかの重要な考えを読み取ることができます。ひとつは誤差についての考え方です。このモデルを$X_j = t_j + e_j$のように，ある個人の変化しない属性について，$j$回測定したとします。この時，測定の平均は次のように計算できます\footnote{この式は，ある測定を多くの個人$i$に対して行ったものとして，$X_i = t_i + e_i$と考えることもできますが，添字が異なるだけで式の展開に違いはありません。}。

\begin{align*}
     & \bar{X}  =  \frac{1}{n}\sum_{j=1}^n X_j                           &             & \mbox{} \\
     & =  \frac{1}{n}\sum_{i=j}^n (t_j + e_j)                            & \mbox{定義より}           \\
     & =  \frac{1}{n} \sum_{j=1}^n  t_j +  \frac{1}{n} \sum_{j=1}^n  e_j & \mbox{分配して}           \\
     & =  \bar{t} + \bar{e}                                              & \mbox{}               \\
\end{align*}

さて古典的テスト理論では，誤差に関して次のことが仮定されます。
\begin{itemize}
    \item 誤差の平均はゼロ。つまり誤差が出現するときは，「正に偏る」「負に偏る」といった一貫した傾向がないと考える。
    \item 真のスコアと誤差との相関はゼロ。つまり誤差は真のスコアに関係なく影響してくるもので，真のスコアと共変動するようであれば偶然によるものとはいえない。
    \item 異なる測定誤差間の相関はゼロ。誤差同士がなにか意味のある変動をしているのであれば，それはもう偶然によるものとはいえない。
\end{itemize}

これらはいずれも，誤差が「偶然によって現れる影響で，制御不可能なもの」という考え方からは自然な仮定だといえるでしょう。より詳しくいえば，誤差は測定に応じて毎回一定の傾向で生じる\keyterm{系統誤差}{けいとうごさ}{systematic error}と，全く傾向のつかめない\keyterm{偶然誤差}{ぐうぜんごさ}{random error}に分けて考えられますが，系統誤差は測定に際して工夫して取り除くべき問題であり，ここでは全くの偶然による誤差についての議論だからです。

さて誤差の平均がゼロ，つまり$\bar{e}=0$ですから，$\bar{X}=\bar{t}$となって，いつかは誤差がなくなって真のスコアを得ることができるようになる，ということが示されます。

また平均は$0$ですが分散はゼロではありません\footnote{ガウスの考えた誤差論から，誤差は確率\keyterm{正規分布}{せいきぶんぷ}{Gaussian Curve}に従うと考えられます。}。このテストの分散を考えると，次のようなことがわかります。
\begin{align*}
     & Var(X)  =  \frac{1}{n}\sum_{i=j}^n (X_j - \bar{X})^2                                                                                              & \mbox{定義より}         \\
     & = \frac{1}{n}\sum_{j=1}^n  ((t_j + e_j) - (\bar{t} + \bar{e}))^2                                                                                  & \mbox{Xをテスト理論のモデルに} \\
     & = \frac{1}{n}\sum_{j=1}^n (  (t_j - \bar{t}) + (e_j - \bar{e}) )^2                                                                                & \mbox{同じ記号でまとめる}    \\
     & = \frac{1}{n}\sum_{j=1}^n (  (t_j - \bar{t})^2 + 2(t_j - \bar{t})(e_j - \bar{e}) + (e_j - \bar{e})^2  )                                           & \mbox{展開する}         \\
     & = \frac{1}{n}\sum_{i=1}^n (t_j - \bar{t})^2 +\frac{1}{n}\sum_{j=1}^n 2(t_j - \bar{t})(e_j - \bar{e})  + \frac{1}{n}\sum_{j=1}^n (e_j - \bar{e})^2 & \mbox{$\sum$を分配}    \\
     & =  Var(t) + 2Cov(te) + Var(e)
\end{align*}
ここで$Cov$とは共分散を表しています。第二項の$2Cov(te)$は真のスコアと誤差との共分散(を2倍したもの)を表していることになりますが，共分散(相関)がそもそも線形関係を表す指標であったことを思い出してください。相関係数は共分散を標準化したものだったわけですが，そういう意味ではこの$Cov(te)$というのは真のスコアと誤差との相関関係を表しているようなものです。さて，ここでも誤差の仮定から，相関はゼロです。すなわち，誤差とはどのような傾向もなく出現するもの，という考えられているのです。どのような傾向もないわけですから，当然なにかと相関関係にあるはずがない，すなわち$Cov(te)=0$であるとなります。

これを踏まえると，$Var(X)=Var(t)+Var(e)$のように，テストの分散が真のスコアの分散と誤差の分散に完全に分割されました。ここから，信頼性の定義は
$$\frac{Var(t)}{Var(X)}=\frac{Var(t)}{Var(t)+Var(e)}$$
と表すことができます。言葉で言えば，\keytermL{信頼性}{しんらいせい}の定義は「全分散中にしめる真のスコアの分散の割合」ということになります。

\section{尺度を評価する}
テストは目に見えないものを測定するためのツールです。できあがったテストがきちんと測定できているのかについては，信頼性と妥当性という2つの側面から評価しなければなりません。
\subsection{信頼性}
信頼性は測定の安定性と言い換えてもいいかもしれません。すなわち，同じものを2回測れば，同じ数字がえられることが安定した測定です。測定するたびに数字が変わるようでは，その測定器(ここでは尺度ですが)は信用ならない，というわけです。
信頼性$Rel$は，次のように表現できるのでした。

\[Rel = \frac{Var(t)}{Var(X)}= \frac{Var(X)-Var(e)}{Var(X)} = 1-\frac{Var(e)}{Var(X)}\]

言葉で表せば，信頼性は全分散に占める真分散の割合であり，全体から誤差分散の割合を引き算したものであるとも言えます。

信頼性のない尺度があれば，その後のお話は先に進みませんから，まずもって「尺度が信頼できるかどうか」を評価する必要があります。このことを\underline{信頼性は妥当性の上限}，と表現したりします。

信頼性を評価する方法は，測定値の安定の程度を評価できればいいのですから，複数の測定を持ってその相関係数を計算することでひとまず達成できます。同じ検査を2回行って，測定値の相関係数を見ることでその安定性を評価する方法を，\keyterm{再検査信頼性}{さいけんさいしんらいせい}{Test-Retest reliability}といいます。

しかし同じ尺度を何度も使うのは，調査回答者に要らぬ構えを持たせてしまいますね。昨日と同じテストを今日もやる，ということになれば対策を練ることができますし，同じ文言が繰り返されると飽きてきてわざと違う答えを書いてやろう，と思われるかもしれません。そこでまったく同じテストではなく，よく似た別のテストを使って一貫性の指標とすることが考えられます。これは\keyterm{平行検査信頼性}{へいこうけんさしんらいせい}{Parallel test reliability}と呼ばれます。

平行検査信頼性は，テストの質が同じぐらいで内容が違うもの，ということになります。算数のテストであれば数字を少し変えることで，同質の異なるものをつくることができたといえるかもしれませんが，他の教科や心理学一般であれば，なかなか難しいことかもしれません。そこで1つのテスト項目を，たとえば偶数項目と奇数項目のように無作為に2分割し，それらの相関係数でもって安定性の測度とすることが考えられます。これを\keytermL{折半法}{せっぱんほう}といいます。

折半法では1つのテスト項目群を半分に分けるのでした。これはなにも半分ではなく，3分割，4分割と細かく割ってそれぞれのブロックの相関係数を計算し，その平均でもって全体的な安定の指標と考えることができます。いやそこまでいくのであれば，5,6,7分割...といって究極的には$M$項目の質問があれば$M$分割してしまえ，ということが考えられますね。

個々の項目が他の項目とどれぐらい相関しているかを考えてみましょう。まず尺度全体から計算される回答者の値は，項目の尺度値の合計であるのが普通です(テストの点数も正答数に対応していますね)。ですから，ある項目$j$の尺度値は，$j$を除いた残り$M-j$個の尺度値の和と高い相関をするはずです。このように，各項目が尺度全体の値とどの程度相関しているかを見る\keyterm{IT相関}{あいてぃーそうかん}{Item-Total correlations}は，尺度の信頼性を見る1つの指標になります。ある項目が，尺度全体と相関していなければ，それはその項目が尺度の中で目的と違うものを測定している可能性があり，それは必然的に尺度の安定を損ねる結果になるからです。

この考え方を発展させ，各項目が他の項目とどの程度相関しあっているか，つまり尺度の中でどの程度整合性がとれたもの，すなわち一貫して同じものを測定しているのかを評価する指標として，\keyterm{$\alpha$係数}{あるふぁけいすう}{alpha coefficient}があります\footnote{\keyterm{クロンバックのアルファ}{くろんばっくのあるふぁ}{Cronbach's alpha}とも呼ばれます。}。これは，テストに含まれる項目数をM，テスト全体の分散を$V_t$，項目jの分散を$V_j$と表すと，次の式で表されます。

\[  \alpha =   \frac{M}{M-1} \times \left(1-\frac{\sum V_j}{V_t} \right) \]

この指標は，各項目が他の項目とどれぐらい相関するかを総合的に表した指標で，とくに\keytermL{内的整合性信頼性}{ないてきせいごうせいしんらいせい}と呼ばれます。

このようにして，尺度の安定の程度である信頼性は数値化して評価されます。安定した尺度であることが確認できた上で，では安定して「何を測っていたのか」を考えることにしましょう。それが妥当性の概念です。

\subsection{妥当性}
\keyterm{妥当性}{だとうせい}{validity}は，信頼性をその上限とした上で，それが何を測っているのかを改めて考える指標です。

たとえば身長を測ろうとして，体重計を使うとします。成長に応じて，身長が伸びますが，それは体重とも関係がありますので，身長の伸びに応じて体重も増えていくでしょう。体重計は安定した計測器で，信頼性は十分あると思いますが，体重計で身長が測れていると言えるでしょうか。相関する変数ですので，部分的にYesといえそうですが，やはり身長は身長計で測ったほうが良いでしょう。身長計の方が，身長という特徴を的確に捉え，より本質に近い測定をしているからです。このように，作ったものがしっかりとその本質を捉えているかどうか，これが妥当性の基本的なポイントです。

妥当性はですから，そもそも概念としてその測定しようとしているものが適切かどうか(\keyterm{構成概念妥当性}{こうせいがいねんだとうせい}{construct validity})とか，その文言でちゃんと質問できているか(\keyterm{内容的妥当性}{ないようてきだとうせい}{content validity})，理屈通りその測定値が結果と変動しているか(\keyterm{基準関連妥当性}{きじゅんかんれんだとうせい}{criterion validity})といった面から検証されます。最後の基準関連妥当性については，基準値と尺度値をつかって数量的に検証できますが，構成概念妥当性や内容的妥当性は中身の問題であったり，言葉と概念の対応であったりするので，数理モデル的アプローチができるものではありません。数値化できないか大きな問題ではない，というのはもちろん逆で，数値化できないところであるからこそ，専門的な観点，幅広い視野，批判的思考でもって検証していかなければなりません。


\section{テスト理論の展開}
\label{NewTestTheory}
ここまでのテスト理論は，古典的テスト理論という名称でした。古典的，というのであれば現代的なテスト理論は何がどうなっているのでしょうか。

現代的なテスト理論，新しいテスト理論ともいわれますが，それは\keyterm{項目反応理論}{こうもくはんのうりろん}{Item Response Theory}，あるいは項目応答理論とよばれます。略して\keytermL{IRT}{IRT}と表現されます。
これを考えるに当たって，テストの項目を分析するというのはどのようなことがなされているのかをみていきたいと思います。

\subsection{通過率と累積正規分布}
みなさんは大学入学共通テスト(旧センター試験，さらにその前は共通一次試験といいました)や，学内の定期テスト，模試など色々なシーンでテストを受けてきたことと思います。模試などでは偏差値が明らかになり，自分の実力が相対的にどのあたりにあるのかがわかるようになっていたかと思います。大学共通テストなどは50万人ぐらいが一度に受験しますから，さまざまな学力の人がそこには含まれるのですが，成績を図にするととても綺麗な正規分布になることが知られています\footnote{\citet{Yamauchi201001}の見返し(表紙を開いた最初の内側のページ)に，センター試験の成績分布が載っており，綺麗な正規分布であることが示されています。}。正規分布は誤差の分布でもありますが，多くの要因が考えられる際の集積的データも，自然とこの形になります。

さて，学力のような潜在変数が標準正規分布に従うと仮定しましょう。この分布の形はどこの確率点がどれぐらいの確率密度を持っているか，あるいはある確率点以上・以下の面積が全体の何\%かを表すものですが，縦軸をある点以下の累積確率に書き直してみましょう(図\ref{fig::04_01})。図\ref{fig::04_01}の上の図がいわゆる正規分布の分布の形，確率密度関数です。下の図はこれを累積確率に書き換えたものになっています。累積確率は0\%から始まって，最終的に100\%にまでどのように増えていくかを示した図になります。
\begin{figure}[ht]
    \begin{center}
        \includegraphics[width=0.8\linewidth,keepaspectratio,keepaspectratio,clip]{../images/Rplot04_01.png}
        \caption{正規分布の確率密度関数(上)と累積確率関数(下)}
        \label{fig::04_01}
    \end{center}
\end{figure}

累積正規確率は，テスト理論と密接な関係があります。というのも，学力が正規分布すると考えるなら，累積正規分布の形はあるテスト項目の\keyterm{通過率}{つうかりつ}{pass ratio}と同じ形になると考えられるからです。

通過率とは，あるテスト項目に正答する人の割合のことです。ここで複数の項目からなる，あるテストをしたとしましょう。正答数を数えてその人の成績とすると，よくできたテストであれば成績は正規分布に従います。さらに，ある項目と成績との相関(\keyterm{IT相関}{ITそうかん}{Item-Total correlations})は高いはずですね。つまりその項目に正答することが，テスト全体の成績と高く関係しているので，その項目はテスト全体が測ろうとしているものを反映していると考えられるからです。また，成績をもとに被験者集団を5群に分けたとしましょう。「成績上位群(HH)」，「成績やや上位(MH)」，「成績中程度(M)」，「成績やや下位(ML)」，「成績下位(LL)」です。このとき，各群の平均通過率を考えると，図\ref{fig::04_02}左上図のようになるのが理想的です。つまり，成績が高い人たちの通過率は高く，成績が低い人達の通過率は低くなるはずですね。同じ図の右上は，LL群でも半分ぐらいが通過し，その後の群は過半数，ほとんどが通過するようになっています。これは，この項目が簡単すぎたことを意味しています。簡単すぎる問題は，それはそれで被験者の特徴が弁別できないという意味で悪い項目です。逆に図の左下にあるのは，HH群でも半分以下の通過率しかありません。つまり難しすぎる問題です。ほとんどの人が間違えてしまうわけですから，これも良い試験問題とは言えないでしょう。右下に至っては逆転していて，どうやったらこういう項目が作れるのか却ってわからないほどですが，学力の低い人だけが正答できて，学力の高い人は正答しない項目ということになります。もちろんこんな項目はよくないわけで，IT相関で負の相関が出ているわけですから，テストの文脈でいうなら「そのテストで測っていない何か別の能力を測っている」と考えるしかありません。

\begin{figure}[ht]
    \begin{center}
        \includegraphics[width=0.8\linewidth,keepaspectratio,clip]{../images/Rplot04_02.png}
        \caption{群ごとの平均通過率。左上が良いパターン。右上は簡単すぎる，左下は難しすぎる項目。右下は逆転していて良くない項目。}
        \label{fig::04_02}
    \end{center}
\end{figure}

ともあれ，このようなやり方で項目の良し悪しを見ていくことができます。また，図\ref{fig::04_02}は5段階ですが，7段階，9段階とよりきめ細かくしていくと，理想的な形は累積正規分布により近づいていきます。新しいテスト理論による項目分析はこの累積正規分布の形を基本とし，これを拡張することで各項目の特徴を描いていくことになります。

ところで1つ前の図\ref{fig::04_01}の下の図には，実線と点線の2つの線が絡んでいることにお気づきでしょうか。実線の方は確率分布関数から累積確率を出して描いたものですが\footnote{Rでは\texttt{pnorm}関数を使って描きます。数式でいうなら，$\int_{-\infty}^p \frac{1}{\sqrt{2\pi}}exp(-\frac{x^2}{2})$となります。}，点線のほうは次の関数を使って描いています\footnote{\texttt{exp}というのは指数関数で，$exp(x)= e^x$のことです。ここで$e$は数学定数で，$e=2.718282...$という数字です。}。

\[ f(x) = \frac{1}{1+exp(-1.7x)} \]

この関数，図から明らかなように累積正規分布とほとんど同じですよね。累積正規分布の関数を直接使うと，積分計算($\int$を使うやつ)が入ってくるのでちょっと計算が面倒ですから，こちらの関数の方を近似関数として用います。この関数のことを\keyterm{ロジスティック関数}{ろじすてぃっくかんすう}{logistic function}と言います。ロジスティック関数そのものは，先の式から1.7という係数を除いた$\frac{1}{1+exp(-x)}$で表されるもので，$-\infty$から$+\infty$までのどんな数字が与えられても，答えを0から1の範囲に変換してしまう関数です。この特徴はとても便利です。というのも，結果が0と1の間に入るということは，比率を表していると考えられるからです。0/1のバイナリデータが従属変数のときに，独立変数をこのロジスティック関数で変換してやれば0か1のどちらに近いか，どれぐらいの比率で1の目が出るかがわかります。項目反応理論も結果が0/1(誤答/正答)ですから，「学力」のような目に見えない能力をロジスティック変換してやれば，結果が正答率になるというのは大変便利なのですね。

それではこのロジスティック関数をつかった項目分析の話に進んでいきましょう。

\subsection{項目母数の特徴}
ロジスティック曲線が累積正規分布の近似関数になっていること，テスト項目の分析には通過率を使って考えることを見てきました。
とくに通過率の分析(図\ref{fig::04_02})では，その項目が難しい設問だったのか，簡単なものだったのかを見ることができました。ロジスティック曲線もこの「項目の難しさ」を表現できるように，次のように拡張できます。

\[ p(\theta) = \frac{1}{1+exp(-1.7 (\theta - b))} \]

左辺の$p(\theta)$に含まれる$\theta$は，潜在特性であり，ここでは標準化された学力のことを想定しますから，偏差値のようなものだと思ってください\footnote{偏差値は標準化スコア$z_i$を$10z_i+50$と変換したものを指します。ここはその変換前の$z_i$と同じです。}。$p(\theta)$は能力$\theta$の人がこの項目に正答する確率=(通過率)です。

ここで$b$という変数が入ってきました。これが\keyterm{困難度}{こんなんど}{difficulty}を表す指標です。$b=0$のときは標準正規分布と同じ形になりますが，$b=1$ならばこの式は右に，$b=-1$ならば左に動きます。つまり$b>0$であれば難しく，$b<0$であれば簡単であることを表現していることになります。図\ref{fig::04_03}に困難度が$b=-1,0,+1$の時の曲線を書いてみましたので確認してください。
\begin{figure}[ht]
    \begin{center}
        \includegraphics[width=0.8\linewidth,keepaspectratio,clip]{../images/Rplot04_03.png}
        \caption{困難度母数の入ったロジスティック曲線。1PLロジスティックモデルともいう。点線が$b=0$の標準的曲線。赤が$b=-1$，青が$b=+1$の例}
        \label{fig::04_03}
    \end{center}
\end{figure}
このように困難度を表現するパラメータを追加したモデルを\keyterm{1パラメータ・ロジスティックモデル}{1ぱらめーたろじすてぃっくもでる}{One Parameter Logistic model}といいます。実際のテストの回答パターンにたいし，各項目にこの曲線を当てはめて困難度を推定することで，項目を評価できるようになります。このように項目の特徴を描く曲線のことを\keyterm{項目特性曲線}{こうもくとくせいきょくせん}{Item Characteristic Curve, ICC}といいます。

さらにパラメータを追加して，次のようにすると\keyterm{2パラメータ・ロジスティックモデル}{2ぱらめーたろじすてぃっくもでる}{Two Parameter Logistic model}になります。
\[ p(\theta) = \frac{1}{1+exp(-1.7 a(\theta - b))} \]
ここで$a$という母数(パラメータ)が入ってきました。これは$a=1$だともとのモデルのままなのですが，これが小さくなると曲線が傾き，大きくなると曲線のカーブが強くなります(図\ref{fig::04_04})。曲線が緩やかになると(図\ref{fig::04_03}の赤線)，$\theta$の違いに対して通過率の変化が乏しくなります。言い換えると感度が悪くなるわけです。逆に曲線の立ち上がりが強くなると(図\ref{fig::04_04}の青線)，$\theta$ がある一定のところを超えるかどうかで正答率がグッと変化することになります。つまりこのパラメータは，回答者の能力$\theta$を分類する力の強さを表しているのです。このパラメータのことをとくに\keyterm{識別力}{しきべつりょく}{discriminant}と呼びます。

\begin{figure}[ht]
    \begin{center}
        \includegraphics[width=0.8\linewidth,keepaspectratio,clip]{../images/Rplot04_04.png}
        \caption{識別力母数の入ったロジスティック曲線。2PLロジスティックモデルともいう。点線が$a=1$の標準的曲線。赤が$a=0.5$，青が$a=1.5$の例}
        \label{fig::04_04}
    \end{center}
\end{figure}

一般にはここまで紹介した2PLモデルがよく用いられます。他にも3，4，5つとパラメータが増えたモデルもありますが\footnote{詳しくは\url{http://antlers.rd.dnc.ac.jp/~shojima/exmk/jindex.htm}を参照してください。3つめのパラメータはあて推量母数，4つ目は上方漸近母数，5つ目は非対称母数と呼ばれています。}，もとのロジスティック曲線に特徴を付け足していったもので，曲線をずらしたり，曲げたりしながらもとのデータにうまく当てはまるようにしつつ，その特徴を解釈できるように工夫しています。

いずれにせよ，テストの結果から各項目の特徴を記述します。少し例示したほうがわかりやすいでしょう。図\ref{fig::04_05}にはとある心理統計の授業で行った試験の結果から，2PLロジスティックモデルを当てはめて項目分析をした例です\footnote{私が本務校で行っている統計の授業では，過去の受講生のデータとさまざまな質問をプールしたデータを貯めてあります。このデータに基づき項目の困難度，テストの平均値を事前に設計し，一貫した基準で評価できるように心がけています。}。同じデータの項目母数を表\ref{tbl::04_01}にも示しました。表\ref{tbl::04_01}の数字と図\ref{fig::04_05}の曲線の対応をよく確認してください。

たとえば，項目\texttt{I0022}は困難度が-1.92，識別力が1.30です。困難度がマイナスですので，これはかなり簡単な問題だということになります。具体的には，偏差値50すなわち$\theta=0$の人であっても98.58\%正解するわけですから，ほとんどの人にとって簡単な問題であることがわかります。
ちなみにこの問題，具体的には帰無仮説検定に関する問いで，「「差がない」「偏りがない」といった仮説は何と呼ばれるか」というものでした\footnote{ちなみに答えは「帰無仮説」です。}。

一方，困難度が0近いところの例として項目\texttt{M0605}をあげますが，これが平均的な難易度の質問になっています。困難度母数$b=0.0$であれば偏差値50，すなわち$\theta=0$の人が正答する確率が50\%の質問ということになりますが，
今回の\texttt{M0605}はそれより少し難しいので，偏差値50の人で20.70\%の割合で正答できます。この問いについて偏差値が70($\theta=2$)であれば，92.96\%の確率で正答できることになりますし，
偏差値30($\theta=-2$)であれば0.51\%，つまりほとんど正答は望めないということになります。ちなみにこの問題は「重回帰分析において，標準化されたデータを使って分析をすることで，モデルの適合度を上げることができる」をYesかNoかで判断させるという質問でした。

他にも項目\texttt{I0017}は困難度が2.17，識別力は0.69です。困難度がもっとも高いグラフで，図の曲線は一番右にあるラインになっています。ただ識別力がやや低いので，曲線はよりなだらかになっていますね。困難度が高いので，$\theta=0$でも7.24\%しか正答できません。難しい！$\theta=2$で45.10\%ですから，かなり能力の高い人でも半分は間違えるような問題です。ちなみにこれは標本分散の期待値が母分散からどれぐらいずれるのかを計算する問題でした。

\begin{figure}[ht]
    \begin{center}
        \includegraphics[width=0.8\linewidth,keepaspectratio,clip]{../images/Rplot04_05.png}
        \caption{実際のテストに2PLモデルを適用した例}
        \label{fig::04_05}
    \end{center}
\end{figure}

% latex table generated in R 4.0.4 by xtable 1.8-4 package
% Wed Mar 17 09:22:06 2021
\begin{table}[ht]
    \centering
    \caption{各項目の困難度と識別力}
    \label{tbl::04_01}
    \begin{tabular}{rlrr}
        \hline
          & 項目ID  & 困難度   & 識別力  \\
        \hline
        1 & I0013 & -1.02 & 1.05 \\
        2 & I0017 & 2.17  & 0.69 \\
        3 & I0022 & -1.92 & 1.30 \\
        4 & I0041 & -0.79 & 0.62 \\
        5 & M0605 & 0.68  & 1.15 \\
        \hline
    \end{tabular}
\end{table}

\subsection{被験者母数の推定}

項目反応理論における潜在特性の推定は，項目の特徴を表す項目母数に対して\keytermL{被験者母数}{ひけんしゃぼすう}と呼ばれ，上述の項目特性に基づいて行われます。
先ほどの図\ref{fig::04_05}をもとに説明します。ある回答者が項目\texttt{I0013}に正答したとしましょう。その人の能力値$\theta$はどの辺りにあるかといえば，確率の曲線に沿った下の領域のどこかということになります(図\ref{fig::04_06})。この曲線，ICCは項目の特徴を表したもので，ある$\theta$の値の能力があればどの程度の確率で正答できるかを表した項目の特徴でもありますが，逆にある人の$\theta$がどのあたりにありそうかを示しているともいえます。
たとえばこのICCにおいて，$\theta$が0.5のときの通過率は86.05\%ですが，言い換えればこの項目に正解した人が$\theta=0.5$である可能性も高そうです。$\theta=-2$の通過率は14.71\%ぐらいですから，能力がこんなに低いとは思えませんし，1つの項目の話でしかないですが，希望的観測をするなら$\theta$がもっと高い可能性もあるでしょう\footnote{$\theta$のありそうな「確率」とはいってないことに注意してください。すぐにわかることですが，このICCの下の面積を積分しても1.0にはなりませんのでこれは確率ではなく，\keyterm{尤度}{ゆうど}{likelihood}のほうなのです！}。

\begin{figure}[ht]
    \begin{center}
        \includegraphics[width=0.8\linewidth,keepaspectratio,clip]{../images/Rplot04_06.png}
        \caption{実際のテスト項目\texttt{I0013}に正答した人の能力がありそうな領域}
        \label{fig::04_06}
    \end{center}
\end{figure}

次に，困難度のより高い項目である\texttt{I0017}には\textbf{誤答}したとしましょう。その人は，\texttt{I0017}のICCの下の領域には\textbf{ない}はずです。項目\texttt{I0013}のICCの下で，かつ，項目\texttt{I0017}のICCの上にあるはずなので，図\ref{fig::04_07}のように塗りつぶされた領域の中に入ります。
\begin{figure}[ht]
    \begin{center}
        \includegraphics[width=0.8\linewidth,keepaspectratio,clip]{../images/Rplot04_07.png}
        \caption{つぎのテスト項目\texttt{I0013}には誤答した人の能力がありそうな領域}
        \label{fig::04_07}
    \end{center}
\end{figure}

この2つの曲線の間にある，濃く彩られた領域の幅が，回答者の能力値がありそうな程度を表しているのです。図\ref{fig::04_07}には$\theta=0$の可能性の大きさを矢印で示してありますが，$\theta$の値はここだけに限らずこの曲線の幅のどこかです。ただ$\theta=2$や$\theta=-2$よりは$\theta=0$のほうが，より「ありそう」な値だということがわかります。

ここでさらに同じ人に，項目\texttt{M0605}の質問をして，この人がそれにも正解したとしましょう。この人の能力$\theta$のありそうな範囲はさらに絞り込むことができます(図\ref{fig::04_08})。
\begin{figure}[ht]
    \begin{center}
        \includegraphics[width=0.8\linewidth,keepaspectratio,clip]{../images/Rplot04_08.png}
        \caption{さらにテスト項目\texttt{M0605}には正答した人の能力がありそうな領域}
        \label{fig::04_08}
    \end{center}
\end{figure}
項目\texttt{I0013}よりは難しい質問に正解したわけですから，$\theta=0$の可能性はグッと小さくなり，それよりも$\theta=2$ぐらいの方がありそうだ，ということになってきます。このように，1つ1つの項目からこの人の能力$\theta$がどのあたりにありそうか，というのを絞っていき，テストに含まれる項目全部を使えば，かなり狭い領域で「この辺りにあるはずだ」と推定できるでしょう。

このIRTを用いた推定方法は，このようにテストの項目ごとにその特徴を分析できるのが特徴です。テストの中でも良い項目，悪い項目というのはあるでしょうが，どのように悪いのかを困難度や識別力といった項目の特徴を使って表現できます。これらの数字は，テストに含まれる項目ごとの難易度を相対的に比較していく中で作られるものであり，回答者の能力に依存するものではありません。テスト理論が被験者の特徴と項目の特徴を分離するところから始まったことを，改めて思い出してください。

また，このような項目同士の比較から求められた項目の特徴をつかって，被験者の能力(潜在特性)を推定する方法についても紹介しました。
その過程の中で気づいたと思いますが，複数の回答を通じてある回答者の能力$\theta$のありそうな領域が狭められていく中で，明らかにその人にとって簡単すぎる問題，難しすぎる問題は意味を成しません。たとえば図\ref{fig::04_08}の段階まで絞り込まれているときに，項目\texttt{I0013}よりも簡単な\texttt{I0022}を出題しても，おそらくほぼ確実に正答し，そのことは「領域を狭める」ことにはなんの貢献もしないでしょう。回答者の能力に相応しい質問を選んで出すことができれば，とても効率よくその領域を絞り込んでいくことができます。こうした考え方は，\keyterm{コンピュータ適応型テスト}{こんぴゅーたてきおうがたてすと}{Computer Adapted Test}として実装されます。

ここではテスト理論について紹介してきましたが，この考え方は後ほど性格テストなどで用いられているリッカート形式の尺度に応用できるように発展します。


\section{現代テスト理論の特徴}
前回は現代テスト理論として項目反応理論(IRT)を紹介しました。ロジスティック曲線を応用して項目の特徴を描画し，それを使って被験者母数を推定する方法についても解説しました。この一連の手続きに基づき，現代テスト理論の利点を考えてみたいと思います。
\subsection{現代テスト理論の利点1:項目母数と被験者母数の分離}
古典的テスト理論からの発展として，現代テスト理論では被験者母数$\theta_i$と項目母数$a_j, b_j$を区分して考えるようになりました。項目母数は通過率のアイデアを精緻にしたものですが，この通過率は項目群の総和を元に考えられていたことを思い出してください。
すなわち，項目母数の計算には項目の相対的な困難度だけを用いています。イメージとしては鉱物の硬度検査のようなものです。2つの異なる硬さの物をぶつけて崩れた方が負け＝より硬度が低いと考えるように，2つの異なる項目を被験者に与えて，より正答者数が少ない方がより難しいと考えるのです。これはつまり，回答者の学力が高かろうが低かろうが，困難度が$b_x < b_y$であるという関係に違いはないという考え方です。

これまでのテストや心理尺度の作り方に比べると，この点が大きく違います。学校などのテストは教員が作っていますから，教員が自分の感覚で「こちらの方がより発展的な内容だ」「こちらの方が難しいだろう」という問いに大きな配点がなされたりするでしょう。その後テストの平均点をみて「今回のテストは簡単にしすぎたか」という判断をしたりするでしょう。しかしIRTでは項目それ自体に困難度を決めさせますから，そこに作成者や回答者の意図は含まれません。平均正答率が高いからといって簡単な問題なのではなく，項目の特徴として困難度が決まるのです。

たとえばサーストン尺度の作り方を思い出してください(セクション\ref{ThurstoneScale},Pp.\pageref{ThurstoneScale}参照)。サーストン尺度では尺度適用前に評定者集団によって尺度値を決めます。この評定者集団が偏った思想の持ち主だけで固められていた場合，尺度の点数は極端なものになり，普通の人がその尺度に回答すると極めて低い点数，高い点数になってしまうかもしれません。
あるいはリッカート尺度の作り方を思い出してください。(セクション\ref{LikertScale},Pp.\pageref{LikertScale}参照)。リッカート尺度では回答者の累積度数から尺度値を算出します。先ほど同様，回答者集団の態度に偏りがあれば，尺度の点数は標準的な物ではなくなるでしょう。
つまり「誰を対象に測定するか」によって目盛りが変わるようなものです。これでは測定結果の一般化は難しいでしょう。たとえばある大学で作られた尺度を，他大学でやってみると違った尺度値になるのですから，研究結果はせいぜい「その大学ではそうなんだろう」となり一般化できなくなります。
従来の方法は，回答者と項目の特徴が関連しすぎていたのです。

これに対し，IRTを使って項目の特徴を計算する場合は，相対的な難易度に違いはありませんから，どの大学で作った尺度であっても統一的な解釈が可能です。尺度作成時に幅広くデータを集め，項目母数を確定させてしまえばどこででも統一的な評価ができます。テストなど学力を測定する際に大学間での違いが見られたとしても，その難易度を調整するのも簡単で，共通する項目を入れておけばそこを基準に相対的な困難度調整ができます\footnote{こうしたテスト間の困難度調整のことをテストの\keyterm{等価}{とうか}{equation}といいます。}。良問と悪問の評価と，回答者の評価を分けることは重要なポイントなのです。
\subsection{現代テスト理論の利点2:被験者母数の推定の利点}
被験者母数と項目母数の分離は，さらに別の利点も生み出します。それはデータが一部欠落した場合の補完に関係します。

リッカート尺度では回答者全体の相対頻度から，カテゴリの尺度値を決定するのでした。ここである人が特定の項目にだけ回答をし忘れたとします(調査研究ではよくあることで，同じような目盛りが並んでいると一行飛ばして丸をつけちゃうようなことはよくあります)。そうすると，その項目だけ合計人数が変わりますから，計算が面倒です。また相関係数を計算する時にも，その人のその項目については計算できなくなります。相関係数を考える時に，一箇所でも欠損値があるとその人のデータを抜いてしまうか\footnote{リストワイズ削除といいます。}，他の値を代入して補完するか\footnote{欠損値補完については，平均値や中央値を代入したり，同じようなパターンで回答している人の値を使い回したり，回帰分析でその項目の値の推定値を入れたり，とさまざまな方法が考えられてきました。欠損値発生のメカニズムにもよりますが，いずれもある程度バイアスのかかった値になってしまいます。統計的によりバイアスの少ない適切な代入法が考えられてはいますが，そもそも欠損値がないのが最も望ましいことに変わりはありません。}，手間でも計算の時にそこだけ外して計算するか\footnote{ペアワイズ削除といいます。}といった工夫が必要です。相関係数に基づいて指標化するときは，欠損値があればその人について算出できないことに変わりはありません。

それに対して，IRTの被験者母数の推定は，項目母数をつかったICCをもとに一人ずつ絞り込んでいくというものでした。もしある人が回答していないことがあっても，計算ができなくなることはありません。その項目の情報が得られないので絞り込み精度は上がりませんが，ヒントが減っただけで回答できないわけではないのです。このように，得られた情報すべてを使ってその人の被験者母数を推定する方法のことを，\keyterm{完全情報最尤推定}{かんぜんじょうほうさいゆうすいてい}{Full Information Maximum Likelihood}といいます。このようにIRTを使うと，必ずしも全問に回答していなくてもスコアは計算できるということになります。欠損値があるからその人のデータは使えない，ということがないのでいいですね！

もっというと，IRTでは全員が全員同じ問題に回答する必要はありません。たとえば能力値が$\theta_i = -0.3$ぐらいにありそうだ，と絞り込めている段階で，次の問題の困難度母数が$b_j = 2.5$であれば，おそらくほぼ確実にその人は回答できないでしょう\footnote{2PLモデルで$a=1,b=2.5$のとき，$\theta=-0.3$が正解する確率は0.8492\%です。}。
その人にいくら困難度母数の高い質問を繰り返しても，ほとんど誤答がつづくだけで，とくにその人の$\theta$がありそうな領域を狭めるヒントにはなりません。むしろ$b_j=-1$とか$b_j=-0.5$のような簡単な問題を出して\footnote{2PLモデルで$a=1,b=-1$のとき，$\theta=-0.3$が正解する確率は76.674\%，$b=-0.5$のであれば58.419\%です。}，それらに正答できるかどうかを見極め，絞り込んで行った方が効率的です。
紙に印刷されているテスト(Paper Based Test)であれば，印刷された問題は変えようがありませんから，困難度順に問題を並べると，あるところから先はずーっと不正解が続く人が続出します。ずーっと不正解なところの問題はいくら良問でも，その人の能力を測るのには役立ちません。ヒントが増えないからです。であれば，回答者の学力に合った問題を，その都度その都度出題した方がいいですよね。コンピュータを使って回答者に相応しい質問をダイナミックに組み替える，コンピュータに基づいたテスト(Computer Based Test)，別名\keyterm{コンピュータ適応型テスト}{こんぴゅーたてきおうがたてすと}{Computer Adaptive Test}というのがそれです。
CATになると，回答者ごとに出題が変わりますからカンニング対策の必要も無くなって，とても便利になること間違いなしです。

\subsection{現代テスト理論の利点3:信頼性についての考え方}
IRTの考え方からいうと，どんな項目でも何らかの情報を提供してくれるはずです。たとえば$b_j=3$のようなとても難しい項目が合ったとします。これがどれぐらい難しいかというと，偏差値70の人でも15\%しか正解できないレベルです。
ほとんどの人にとっては誤答にしかならない難しすぎる悪問だ，といいたくなるかもしれませんが，学力が高い人がどれほど高いレベルでやれるのかを検証するためには必要な問題です。偏差値が70なのか，75なのか，80までいけるのか，といったことを見極めるためにはこの問題でないと情報が得られないことになります。逆に簡単すぎる問題であっても，より低いレベルを精緻に検証するためには必要なものなのです。

つまり，どの項目にもその項目が得意とする領域があるはずです。この項目はこの領域の回答者を絞り込む時に，最も有用な情報をもたらしてくれるはず，という$\theta$の場所があるはずなのですね。
これを表現するのが\keyterm{項目情報曲線}{こうもくじょうほうきょくせん}{Item Information Curve}といい，次の式で表される項目情報関数で描くことができます。

\[ I(\theta) = a_j^2 p_j(\theta)q_j(\theta)\]

ここで$p_j(\theta)$はその項目$j$の$\theta$における正答率(通過率)，$q_j(\theta)$は誤答率を表しています。能力が平均的，すなわち$\theta=0$のときは，$0.5 \times 0.5$になる平均的な困難度($b=0$)の設問が最も大きな値になる，というわけですね。
図\ref{fig::05_01}にいくつかのICCとそれに対応するIICを描きましたので，確認してください。
\begin{figure}[ht]
    \begin{center}
        \includegraphics[width=0.8\linewidth,keepaspectratio,clip]{../images/Rplot05_01.png}
        \caption{項目特徴曲線(ICC,上図)と項目情報曲線(IIC,下図)の例。左から順に$a_1=0.8,b_1=-1$の識別力が弱く簡単な項目，$a_2=1,b=1$のやや困難な項目，$a_3=1.2,b_3=2$の識別力が高く困難な項目。}
        \label{fig::05_01}
    \end{center}
\end{figure}
IICのピークは，対応するICCの$\theta=0.5$のところにあること，識別力はピークの尖り具合に関わっていることを確認してください。

さて，IICはその項目がどこで情報をもたらしてくれるか，ということを表しているのでした。言い方を変えると，IICのピークはその項目のもっとも信頼できるところであるともいえます。つまりIRTにおいて\textbf{信頼性は潜在特性の関数}になっているのです。
古典的テスト理論ではテスト全体の分散に占める真分散の割合のことを信頼性というのでした。IRTでは項目ごとにこれを考え，さらにその項目のもっとも感度の良いところを探る関数として，その信頼性を評価できるようになったといえるでしょう。

またIICはある項目から得られる情報のことを意味しますが，テストに含まれているすべての情報関数を足し合わせることで，そのテストから得られる情報の大きさを関数として評価できます。すなわちテスト全体の情報量$I_T(\theta)$は，$I_T(\theta)=\sum_{j=1}^M I_j(\theta)$であり，この関数のことを\keyterm{テスト情報関数}{てすとじょうほうかんすう}{Test Infromation Curve}といいます。図\ref{fig::05_01}の3つの項目からなるTICを示したのが図\ref{fig::05_02}です。
\begin{figure}[ht]
    \begin{center}
        \includegraphics[width=0.8\linewidth,keepaspectratio,clip]{../images/Rplot05_02.png}
        \caption{テスト情報関数(黒の実線部)。理解を進めるために各IICを薄い点線で表現した。}
        \label{fig::05_02}
    \end{center}
\end{figure}
これを見ると，この3つの項目からなるテストは$\theta=2$より少し低いレベルを測定するときに最も鋭敏に働くということがわかります。このように，項目母数がわかっていれば事前にテストの特徴をどのあたりに持ってくるかを決定でき，自由自在にテストをデザインできるようになるわけです。

テストの例で話をしていますが，心理学的な領域でももちろん便利な手法です。たとえば高い認知能力レベルの人をとくに選出したいとか，精神的な健康度がごく低い人をしっかりと検出したいといった目的があれば，そのあたりにピークが来るような項目からなる質問項目からなる調査票を構成すれば良いのです\footnote{具体例として\citet{kosugi120005596041}をあげておきます。学校適応感を測定するため，テストのピークがやや低いところに来るようになっています。}。


\subsection{現代テスト理論の問題点}
ここまで見てきたように，IRTはさまざまな利点があります。しかし欠点がないわけではありません。

ここまでの話はすべて，項目母数が既にわかっていれば，という前提付きで進めてきました。では項目母数はどのようにして定めるのでしょう？ これはもちろん得られたデータから算出できるのですが，そのためには事前に多くの被験者から回答を集め，項目母数の値はほぼ間違いなくこれぐらいだろう，といえるほど安定したものである必要があります。テストの場合，回答が0/1というバイナリデータで得られますから，そもそもそれほど情報がある反応ではありません。バイナリデータからその項目の特徴を安定して推定するためには，かなり多くの被験者(数千から数万単位)を集めて項目に回答させておく必要があります。テスト項目は一度使ってみないと，項目母数がどうなるかわからないというのもポイントです。

また，CATなど項目をダイナミックに組み合わせるためには，選べるぐらいさまざまな項目を準備しておかなければなりません。項目を集めたものを\keyterm{項目プール}{こうもくぷーる}{Pool of Items}といいますが，これも数千から数万の単位で用意しておく必要があります。なぜなら，テスト項目は事前に1回は使ってあるわけですから，数えるほどしか項目がなければ受験生が正答を事前に丸暗記できてしまうからです。もちろん項目プールが数千から数万あっても一度どこかで使われていますから，過去問をすべて完全に丸暗記すればその人は満点が取れてしまいます。もっともそれだけのものを覚えられるのは，ある意味学力が高いといっても差し支えないと思いますが。

日本でおこなわれる大学入学共通試験をはじめ，試験問題というのは極めて厳重な管理下に置かれ，事前にその情報が漏れることは公平性の観点から言って不適切であるとされています。しかしIRTで分析するためには，事前に項目の特徴を知っていなければならないのです。CATをつかって入学試験などができれば，カンニング対策にもなりますし，受験生は何度でもチャレンジできるので利点も多いのですが，「公平性のために新しいテストでなければならない」となるとなかなか実用化できないところがあるというのも事実です\footnote{令和2年度に大学入試センター試験から大学入学共通試験に変わりましたが，改革前の計画ではCAT化することが盛り込まれていました。しかし実際には，受験生のためのコンピュータやタブレットを準備したり，安定した通信網が必要であったり，というハード的な問題もあって見送られてしまいました。}。

ところで，心理学の場合は学力テストのように正答・誤答があるものではなく，「当てはまる」から「当てはまらない」といった軸上で多段階の反応を求めることが一般的です。それでは多段階反応に拡張されたモデルを見ていきましょう。
\section{段階反応モデル}
リッカート法などで作られる尺度は，一般に5，7段階のものが多くあります。少ないものでは3件法\footnote{たとえばYG性格検査は3段階です。}であったり，ものによっては4，6件法であることもあります\footnote{偶数の段階にすることで，必ずどちらかの極に寄るようにして集計できます。日本人は「どちらでもない」に回答しがちな中点集中傾向があるとも言われているので，わざと肯定・否定のどちらかに寄せようという考え方です。}。しかしこれらの段階はいずれも順序尺度水準の情報しか持っておらず，そのままでは尺度値として使うことができません。シグマ法などで数値化すれば良いのですが，その手間を省いて分析する悪い習慣もあることは既に指摘した通りです。

IRTの多段階版はこうした問題に対応できる方法です。IRTの多段階モデルは大きくわけて2つあり，1つは\keyterm{段階反応モデル}{だんかいはんのうもでる}{Graded Response Model; GRM}\citep{samejima1997graded}，もうひとつは\keyterm{多段採点モデル}{ただんさいてんもでる}{Partial Credit Model}\citep{muraki1992generalized}と呼ばれています。どちらも発想は似たようなところがあり，ここではGRMについて解説します。興味がある人は，\citet{Toyoda20120824}や\citet{Katoken2014}など専門書を参考にしてください。

GRMの考え方の基本は，段階反応の背後には正規分布する連続的な潜在特性$\theta$がある，と仮定するところです。心理的な能力，学力，性質などは連続的なのですが，それが表に出てくる時は離散的(順序的)だというわけです。図\ref{fig::05_03}に図示されているように，正規分布がある\keyterm{閾値}{いきち}{threshold}(これを$b_k$と表しますが)を超えると出現する時は次のカテゴリになる，ということを考えます。図は三段階の例ですが，図から明らかなように$k$段階であれば閾値の数は$k-1$個あることになります。横軸$\theta$は心理的な態度や性質の強さだと思ってください。さてそうすると，$\theta=2.0$ぐらいであれば，ほぼ間違いなく「当てはまる」に回答することになりますし，$\theta=-2$であれば「当てはまらない」に回答するようになるはずです。このように$\theta$が大きくなればなるほど最後のカテゴリに反応する確率は上がっていきますから，ここは2PLモデルの時のようにロジスティック曲線で「当てはまる」に回答する確率を表現できるでしょう。問題は，それより下の段階に反応する確率をどのように表現するか，です。

ここである段階に反応する確率を考えるために，少し表現を改めます。すなわち，個人$i$の項目$j$に対する反応が，カテゴリ$k$\textbf{以上}になる確率として，$P^{+}_{jk}(\theta) = P(x_{ij} \ge k|\theta)$をまず考えます。ここで$k=0,1,2,...,K$とします。先ほど示したように，$k=K$，すなわち一番上のカテゴリ(ここでは「当てはまる」)に回答する確率は，2PLロジスティック関数と同じ形をしていますから，次のように表現できます。

\[ P_{jK}(\theta) = P^{+}_{jK}(\theta) = \frac{1}{1+exp(-a_j(\theta-b_{jK}))} \]

ここで右辺の$a_j$は\keytermL{識別力}{しきべつりょく}，$b_{jK}$はカテゴリ$K$の\keytermL{困難度}{こんなんど}を表しています。
左辺の$P_{jK}(\theta)$は$\theta$の人が項目$j$のカテゴリ$K$に反応する確率で，それは$k$以上に反応する確率$P^{+}_{jk}(\theta)$と一致していることを表しています。

次に，もっとも低い段階に回答する確率を考えましょう。これは$\theta$が大きくなればなるほど減っていくはずで，いわばロジスティック曲線の逆のような形をするはずです。反応確率は最大でも1.0ですから，逆ということは1.0から引いてやればよいでしょう。

\[ P^{+}_{j0}(\theta) = 1.0 -   \frac{1}{1+exp(-a_j(\theta-b_{j0}))} \]

問題は「どちらでもない」に反応する確率です。これは引き算で考えることができます。すなわち「どちらでもない」以上に反応する確率から，「当てはまる」以上に反応する確率を引いてやれば良いのです。

\[P_{jk}(\theta) = P^{+}_{jk}(\theta)-P^{+}_{j,k+1}(\theta) \]

ここにあるように，段階数が増えたとしても$k$番目のカテゴリ以上に反応する確率から，$k+1$番目に反応する確率を引いてやれば，$k$番目のカテゴリに反応する確率が得られます。この計算をして描かれる曲線のことを\keyterm{項目反応カテゴリ特性曲線}{こうもくはんのうかてごりとくせいきょくせん}{Item Response Category Charactaristic Curve; IRCCC}，あるいは単に\keyterm{カテゴリ確率曲線}{かてごりかくりつきょくせん}{Category Probability Curve}と呼ばれます。IRCCCは図\ref{fig::05_03}の下段に示されています。

\begin{figure}[ht]
    \begin{center}
        \includegraphics[width=0.8\linewidth,keepaspectratio,clip]{../images/Rplot05_03.png}
        \caption{正規分布と閾値(上図)とIRCCC(下図)}
        \label{fig::05_03}
    \end{center}
\end{figure}

IRCCCを，$\theta$の値がマイナスからプラスの方向に動かしながら見ていってください。最初は当然「当てはまらない」に反応する確率が一番高いのですが，それが徐々に下がっていきます。「どちらでもない」の反応確率は徐々に増えていき，閾値$b_{j1}$で「当てはまらない」と逆転しピークを迎えることになります。その頃「当てはまる」も徐々増えていき，閾値$b_{j2}$でピークが逆転する，というようになります。ピークは逆転されても，他の確率が0になっているわけではなく，可能性は残っています。
またIRCCCもICC同様に変換して，情報曲線に帰ることができます。すなわちどのあたりで鋭敏に情報を検出できるかを表現する項目反応カテゴリ情報曲線を描くこともできます。このようにして，段階反応でもその項目の特徴をデザインできるのです。

\subsection{適切な反応段階を考える}
\begin{figure}[ht]
    \begin{center}
        \includegraphics[width=0.8\linewidth,keepaspectratio,clip]{../images/Rplot05_04.png}
        \caption{適切な反応段階をデザインする}
        \label{fig::05_04}
    \end{center}
\end{figure}
実際の調査研究をすると，図\ref{fig::05_04}の上の段のようなIRCCCが描かれることがあります。何かおかしいところがありませんか？
これは5段階の反応モデルですが，4番目の反応カテゴリがずいぶん低く，そのピークが3番目と5番目の反応カテゴリに潰されてしまっていますね。つまり，4番目の反応がもっとも際立つシーンがないということです。言い換えるならば，これは尺度作成側が5段階だと思っていたにもかかわらず，回答者はどういう時に「やや当てはまる」と答えるのかがはっきりせず，実質4段階でしか反応していないことを表しています。

このようなIRCCCが描かれてしまう場合は，$k=3$の反応と$k=4$の反応を合わせてひとつにしてしまうなど，段階の修正を考えると良いでしょう。具体的にはデータで4と入力していたものを，3に置き換えたりします\footnote{4を5に書き換えても構いません。ヒストグラムを見て，より正規分布っぽくなるようにすると良いでしょう}。修正したのが下の図になります。このように修正しても，情報関数は変わりません。同じデータから得られる情報は同じだからです。

このように5段階，7段階を設定して回答者に無理やり回答を求めても，分析するとカテゴリのピークが潰れていることがあります。回答者の反応しやすいカテゴリ数を丁寧に設計してやることが重要です。もちろん無分別に尺度値をつけて，そのまま分析するのはもっとも不適切な方法です。


\end{document}